% Preámbulo para la exportación del informe TI-12 con pandoc + xelatex
\usepackage{graphicx}
\usepackage{xcolor}
\usepackage{booktabs}
\usepackage{longtable}
\usepackage{array}
\usepackage{ragged2e}
\usepackage{fancyhdr}

% Tablas: se reducen y se permite el salto de línea dentro de la celda
\AtBeginEnvironment{longtable}{\footnotesize\setlength{\tabcolsep}{3pt}}
\AtBeginEnvironment{tabular}{\footnotesize\setlength{\tabcolsep}{3pt}}
\renewcommand{\arraystretch}{1.08}

% Bloques de código: los diagramas ASCII necesitan interlineado justo y no partirse
\usepackage{fvextra}
% pandoc define sus propios Verbatim más abajo: se fuerza al inicio del documento
\AtBeginDocument{\fvset{fontsize=\scriptsize,breaklines=false,breakanywhere=false,samepage=false}}

% Encabezado y pie
\pagestyle{fancy}
\fancyhf{}
\fancyhead[L]{\scriptsize TI-12 · Cumplimiento normativo en proyectos TIC}
\fancyhead[R]{\scriptsize audIT · Empresa N.º 10}
\fancyfoot[C]{\thepage}
\renewcommand{\headrulewidth}{0.3pt}

% Títulos algo más compactos
\usepackage{titlesec}
\titlespacing*{\section}{0pt}{1.4ex plus .2ex}{0.8ex}
\titlespacing*{\subsection}{0pt}{1.1ex plus .2ex}{0.6ex}
\titlespacing*{\subsubsection}{0pt}{0.9ex plus .2ex}{0.5ex}

% Sin sangría, con separación entre párrafos
\setlength{\parindent}{0pt}
\setlength{\parskip}{0.55em}
